
@@M0@@
@@M1@@
Et pour établir II, il suffit de voir que, quitte à restreindre @@M2@@, @@M3@@, @@M4@@ et @@M5@@ sont quasi-affines.

Considérons @@M6@@ d'abord. En remplaçant @@M7@@ par @@M8@@ et @@M9@@ par @@M10@@, on se ramène au cas @@M11@@ intègre sur @@M12@@. Comme @@M13@@ est surjectif, il existe l'un des @@M14@@, soit @@M15@@, qui domine @@M16@@, et on a les deux suites
@@M17@@
@@M18@@
Si l'on sait donc que les @@M19@@ donc aussi @@M20@@ sont représentables, et le morphisme composé est quasi-affine, on conclut que @@M21@@ est quasi-affine. Compte tenu de la démonstration de 1.2 on est donc ramené pour @@M22@@ à démontrer II'. \uline{Le théorème} 1.1 \uline{lorsque} @@M23@@ \uline{et} @@M24@@ \uline{sont intègres sur} @@M25@@.

Pour @@M26@@ et @@M27@@, on doit traiter le cas d'un morphisme @@M28@@ où @@M29@@ est réduit, les @@M30@@ sont les composantes irréductibles de @@M31@@ et sont intègres sur @@M32@@. Or on peut supposer @@M33@@ à fibres connexes. En effet, soit @@M34@@ le point géométrique de @@M35@@. Alors @@M36@@ a comme composantes irréductibles les @@M37@@ et, quitte à restreindre @@M38@@, on peut supposer que @@M39@@ si et seulement si
